\section*{Aktivitas Pembelajaran dan Praktikum}
\addcontentsline{toc}{section}{Aktivitas Pembelajaran dan Praktikum}

\begin{table}[H]
\centering
\caption{Rencana aktivitas Bab III}
\label{tab:ch03-aktivitas}
\begin{tabularx}{\textwidth}{@{}llX@{}}
\toprule
\textbf{Sesi} & \textbf{Bentuk} & \textbf{Kegiatan} \\
\midrule
1 & Ceramah & Keluarga translator dan perbedaan compiler-interpreter \\
2 & Diskusi & Fase kompilasi dan aliran informasi \\
3 & Latihan & Membaca tabel simbol dan pesan kesalahan \\
4 & Praktikum & M0: menyiapkan repositori dan spesifikasi Subset-C \\
\bottomrule
\end{tabularx}
\end{table}

\subsection*{M0: Repositori dan Spesifikasi Bahasa}
\textbf{Tujuan.} Menyiapkan struktur proyek dan spesifikasi awal
Subset-C sebagai fondasi seluruh milestone berikutnya.

\textbf{Langkah kerja.}
\begin{enumerate}[label=\arabic*.,leftmargin=*]
  \item Salin direktori \file{code/subsetc} ke repositori pribadi.
  \item Baca \file{code/subsetc/spec/subset-c.md} dan
        \file{code/subsetc/README.md}.
  \item Jalankan \file{build.bat} atau \file{Makefile} untuk memeriksa
        ketersediaan toolchain. Jika toolchain belum ada, catat
        langkah instalasinya.
  \item Tambahkan minimal dua program contoh baru ke
        \file{examples/}: satu program aritmetika dan satu program
        dengan kontrol alur.
  \item Tuliskan satu halaman catatan tentang keputusan desain bahasa
        yang Anda ambil.
\end{enumerate}

\textbf{Definisi selesai.}
\begin{checklist}
  \item Repositori memiliki struktur \file{spec}, \file{include},
        \file{src}, \file{examples}, \file{tests}.
  \item Spesifikasi v0.1 terbaca dan dipahami.
  \item Minimal tiga program contoh tersedia.
  \item Catatan keputusan desain diserahkan.
\end{checklist}

\textbf{Penilaian.} M0 tidak memiliki bobot tersendiri, tetapi menjadi
prasyarat Praktikum 1 (8\%). Rubrik yang dipakai adalah rubrik umum
praktikum pada Lampiran A.
